\begin{textsample}{POS Dim 2 – 2011\_c5 – Score -1.00 – 2011\_c5/enade\_2011\_art\_vis\_14aea27cb0.txt}  \label{ex:f2_pos_348}
agregações resultantes dos seguintes critérios : nível nacional e por Grande Região, Categoria Administrativa e Organização Acadêmica.
O Capítulo 4 trata das percepções dos estudantes quanto à prova ENADE/2011, as quais foram analisadas por meio de nove perguntas que avaliaram desde o grau de dificuldade do exame até o tempo gasto para resolver as questões.
Nesse capítulo objetivou -se a descrição desses resultados, relacionando os estudantes a quatro grupos de desempenho ( limitados pelos percentis : 25 \% ; 50 \% ou mediana ; e 75 \% ), bem como às Grandes Regiões onde os cursos estavam sendo oferecidos.
O Capítulo 5 expõe o panorama nacional da distribuição dos conceitos dos cursos avaliados no ENADE/2011, por meio de tabelas e análises que articulam os conceitos à Categoria Administrativa e à Organização Acadêmica, estratificadas por Grande Região.
O Capítulo 6 enfatiza as características dos estudantes, reveladas a partir dos resultados obtidos no \textbf{Questionário} do Estudante.
O estudo desses dados favorece o conhecimento e a análise do perfil socioeconômico, a percepção sobre o ambiente de ensino-aprendizagem e dos fatores que podem estar relacionados ao desempenho dos estudantes, cujas características são articuladas ao seu desempenho na prova, à Grande Região de funcionamento do curso e à Categoria Administrativa da IES.
Espera -se que as análises e resultados aqui apresentados possam subsidiar redefinições político-pedagógicas aos percursos de formação no cenário da educação superior no país.
CAPÍTULO 5 DISTRIBUIÇÃO DOS CONCEITOS 5.1 PANORAMA NACIONAL DA DISTRIBUIÇÃO DOS CONCEITOS A Tabela 5.1 apresenta a quantidade e distribuição de cursos de Artes Visuais ( Licenciatura ) participantes do ENADE/2011, por faixa de conceito e Grande Região.
A diferença entre os cursos tabulados neste capítulo e no capítulo 2 são os cursos sem conceito, em princípio, aqueles sem alunos concluintes que participassem da prova.
Observando -se os dados da Tabela 5.1, nota -se que, dos 94 cursos participantes, 41 ( 43,6 \% ) classificaram -se com conceito 3, o valor modal.
Este foi também o conceito modal nas regiões Norte ( 60,0 \% ), Sudeste ( 48,7 \% ) e Centro-Oeste ( 50,0 \% ).
Nas regiões Nordeste e Sul o conceito modal foi 4, respectivamente, 55,6 \% e 45,2 \%.
O conceito 4 foi o segundo mais frequente em nível nacional ( 33,0 \%, correspondendo a 31 cursos ) e o conceito 2, o terceiro ( 9,6 \%, correspondendo a nove cursos ).
Houve, ainda, sete cursos ( 7,4 \% ) que receberam conceito 5 e outros três cursos que receberam conceito 1.
Dos 94 cursos de Artes Visuais ( Licenciatura ) três ficaram sem conceito ( SC ).
A região Norte participou com cinco cursos ou 5,3 \% do total nacional.
Nesta região, o conceito 3 foi atribuído a três cursos dos cinco participantes, o que equivale a 60,0 \% do total regional.
Nenhum curso ficou sem conceito.
Os demais foram avaliados com conceito 2 ( dois cursos, correspondendo a 40,0 \% ).
A região Nordeste participou com nove cursos ou 9,6 \% do total nacional.
Como já comentado, destes, dois cursos, 22,2 \% em termos regionais, obtiveram conceito 3.
O conceito 4 foi atribuído a cinco cursos ( 55,6 \% ), o conceito modal para a região.
O conceito 5 foi atribuído a dois cursos ( 22,2 \% ).
Nessa região, nenhum dos cursos ficou sem conceito.
Dos 43 cursos participantes da região Sudeste, 21 ( 48,7 \% ) obtiveram conceito 3, o conceito modal.
O conceito 1 foi atribuído a três cursos ( 7,0 \% ) e o conceito 2, a sete ( 16,3 \% ).
Oito cursos ( 18,6 \% ) receberam o conceito 4 e dois outros ( 4,7 \% ) receberam o conceito 5.
Dois cursos ficaram sem conceito.
A região Sul também contou com cursos distribuídos em quase todas as faixas de conceitos.
A predominância do conceito 4 foi de 45,2 \%, correspondentes a 14 dos 31 cursos participantes na região Sul.
O conceito 3 foi atribuído a 12 cursos ( 38,7 \% ) e o conceito 5, a três cursos ( 9,7 \% ).
O conceito 2 foi recebido por um único curso ( 3,2 \% ) e nenhum curso recebeu o conceito 1.
Um dos cursos da região Sul ficou sem conceito.
A metade ( 50,0 \%, correspondendo a três cursos ) dos seis cursos participantes na região Centro-Oeste recebeu conceito 3, o conceito modal.
Os demais cursos foram avaliados com conceito 2 ( um curso, 16,7 \% ) e conceito 4 ( dois cursos, 33,3 \% ).
Nesta região nenhum curso ficou sem conceito.
Nenhum curso, tampouco, foi alocado nos conceitos 1 ou 5. 5.2 CONCEITOS POR CATEGORIA ADMINISTRATIVA E POR GRANDE REGIÃO A Tabela 5.2 apresenta a distribuição dos cursos participantes do ENADE/2011 de Artes Visuais ( Licenciatura ), por Categoria Administrativa, de acordo com os conceitos por eles alcançados, segundo as Grandes Regiões brasileiras.
Dos 94 cursos participantes, 39 ( 41,5 \% ) eram ministrados em instituições públicas e 55 ( 58,5 \% ), em privadas.
De acordo com as informações da Tabela 5.2, em termos nacionais houve concentração dos cursos nos conceitos mais elevados, com 79 cursos ( 84,0 \% ) obtendo conceito igual ou superior a 3.
Dos 39 cursos participantes de IES públicas, o conceito 4 foi o valor modal, atribuído a 16 cursos.
Entre os cursos das demais IES públicas participantes, um obteve conceito 1 ( 2,6 \% da categoria ), outro recebeu conceito 2 ( 2,6 \% ) e 14 foram alocados ao conceito 3 ( 35,9 \% ).
Os seis cursos restantes ( 15,4 \% ) receberam o conceito 5.
Nesta categoria, um dos cursos ficou sem conceito.
Na rede privada, o conceito modal foi 3, com 27 cursos dos 55 da categoria.
Entre os demais cursos participantes, dois receberam conceito 1 e oito receberam conceito 2.
O conceito 4 foi atribuído a 15 cursos e o conceito 5 a um curso.
Nesta categoria administrativa, dois dos cursos ficaram sem conceito.
Na análise por região, observa -se que na região Norte, as instituições privadas participaram com dois dos cinco cursos ( 40,0 \% do total regional ), dos quais nenhum ficou sem conceito.
O conceito para as instituições privadas na região ficou entre 3 e 4, cada um com um curso.
O conceito para esta Categoria Administrativa, na região Norte, também ficou distribuído entre 3 e 4, com dois e um curso cada, respectivamente.
Nesta categoria e região nenhum curso foi alocado nos demais conceitos.
As instituições públicas participaram com três cursos ( 60,0 \% do total regional ), dos quais nenhum ficou sem conceito.
O conceito modal foi 3, com dois cursos. região nenhum curso foi alocado nos demais conceitos.
As instituições públicas participaram com três cursos ( 60,0 \% do total regional ), dos quais nenhum ficou sem conceito.
O conceito modal foi 3, com dois cursos.
Na região Nordeste, a rede privada contou com apenas um dos nove cursos participantes, equivalentes a 11,1 \% do total da região, a menor proporção entre todas as regiões.
Este único curso obteve conceito 3.
As instituições públicas dessa região participaram com oito dos nove cursos ( 88,9 \% \% ), dos quais cinco obtiveram conceito 4, o valor modal.
Os três demais foram avaliados com os conceitos 3 ( um curso ) e 5 ( dois cursos ).
Nenhum curso ficou sem conceito nesta combinação de região e categoria.
Na região Sudeste, a proporção de cursos da rede privada, 76,7 \%, foi mais elevada do que nas demais regiões brasileiras, correspondendo a 33 dos 43 cursos participantes.
Nesta categoria, na região Sudeste, o conceito modal foi 3 ( 17 cursos ).
Os demais foram avaliados com conceito 1 ( dois cursos ), conceito 2 ( seis cursos ), conceito 4 ( cinco cursos ) e conceito 5 ( um curso ).
Nesta combinação de Categoria Administrativa e Grande Região, dois cursos não receberam conceito.
Entre os dez cursos em instituições públicas na região Sudeste, a categoria modal foi o conceito 3 ( quatro cursos ).
Os demais foram avaliados com os conceitos : 1 ( um curso ), 2 ( um curso ), 4 ( três cursos ) e conceito 5 ( também um curso ).
Na região Sudeste nesta categoria, nenhum curso ficou sem conceito.
As instituições privadas concentraram 16 dos 31 cursos participantes da região Sul, 51,6 \% do total regional.
Desses, oito ficaram com conceito 4, o conceito modal.
Os demais foram avaliados com conceito 2 ( um curso ) e conceito 3 ( sete cursos ).
Nesta combinação de Categoria Administrativa e Grande Região, nenhum curso ficou sem receber conceito ou recebeu o conceito máximo.
As instituições públicas na região Sul participaram com 15 cursos ( 48,4 \% ), dos quais um ficou sem conceito.
Seis cursos ficaram com conceito 4, o conceito modal.
Os demais foram avaliados com conceito 3 ( cinco cursos ) e conceito 5 ( três cursos ).
Nenhum curso foi alocado aos conceitos 1 e 2.
Na região Centro-Oeste, houve um equilíbrio na divisão por Categoria Administrativa com três cursos para cada tipo de IES.
Os cursos de instituições privadas foram avaliados nos conceitos 2, 3 e 4, com um curso cada.
Dos três cursos de instituições públicas, dois foram avaliados no conceito 3, conceito modal, e outro no conceito 4.
Nesta região, nenhum curso ficou sem conceito. 5.3 CONCEITOS POR ORGANIZAÇÃO ACADÊMICA E POR GRANDE REGIÃO Na Tabela 5.3 encontra -se a distribuição dos conceitos atribuídos aos cursos participantes do ENADE/2011 na Área de Artes Visuais ( Licenciatura ), por Organização Acadêmica, segundo as Grandes Regiões brasileiras.
Dos 94 cursos participantes, 61 eram oferecidos em Universidades, 14 em Centros Universitários e os demais 19 em Faculdades.
Esta distribuição corresponde a, respectivamente, 64,9 \%, 14,9 \% e 20,2 \% dos cursos.
De acordo com os dados apresentados, dos sete cursos avaliados com conceito 5 seis eram vinculados a Universidades.
Um curso em Universidades ficou sem conceito.
Esse tipo de Organização Acadêmica teve o conceito 4 como modal, com 26 cursos.
Os demais cursos avaliados receberam os conceitos 1 ( um curso ), 2 ( três cursos ), 3 ( 24 cursos ) e conceito 5 ( seis cursos, como já mencionado ).
Entre os cursos em Centros Universitários, o conceito modal foi 3, com sete cursos.
Neste tipo de organização acadêmica um curso ficou sem conceito e nenhum recebeu o conceito 5.
Os outros cursos neste tipo de Organização Acadêmica receberam os conceitos 1 ( um curso ), 2 ( dois cursos ) e conceito 4 ( três cursos ).
Nas Faculdades, um dos 19 cursos ficou sem conceito e dez receberam o mesmo conceito modal dos Centros Universitários, conceito 3.
Dos demais cursos neste tipo de Organização Acadêmica, um recebeu conceito 1, quatro conceito 2, dois cursos conceito 4 e um curso conceito 5.
Considerando -se separadamente as regiões brasileiras, verifica -se que na região Norte, as Universidades concentraram quatro dos cinco cursos participantes.
Dos cursos em Universidades, estes ficaram distribuídos igualmente nos conceitos modais 3 e 4.
Os Centros Universitários não tiveram representação.
A Organização Acadêmica Faculdade foi representada por um único curso que recebeu o conceito 3.
Na região Nordeste, as Universidades participaram com todos os cursos na Área de Artes Visuais ( Licenciatura ) da região, um total de nove.
Nenhum dos cursos em Universidades no Nordeste ficou sem conceito.
O conceito modal foi 4, com cinco cursos.
Os demais receberam conceitos 3 e 5, com dois cursos cada.
Na região Sudeste, as Universidades concentraram 23 dos 43 cursos da região.
Entre os cursos em Universidades na região, o conceito modal foi 3 com 12 cursos, e nenhum curso ficou sem conceito.
Os demais cursos receberam os conceitos 1 ( um curso ), 2 ( três cursos ), 4 ( cinco cursos ) e 5 ( dois cursos ).
Os Centros Universitários participaram com 11 cursos na região Sudeste, dos quais seis obtiveram conceito modal, 3, e um curso ficou sem conceito.
Os demais receberam os conceitos 1 ( um curso ), 2 ( dois cursos ) e conceito 4 ( um curso ).
Tampouco participaram cursos nesta combinação de organização acadêmica e região com conceito 5.
As Faculdades foram representadas por nove cursos na região Sudeste, que se distribuíram nos conceitos : 3 ( três cursos, conceito modal ), 1 ( um curso ), 2 e 4 ( dois cursos, cada ).
Nenhum curso recebeu conceito 5.
Dos 31 cursos da região Sul, 22 eram de Universidades, para os quais o conceito modal foi 4, com 13 cursos.
Nesse tipo de organização, um dos cursos ficou sem conceito e os demais receberam os conceitos 3 ( seis cursos ) e 5 ( dois cursos ).
Os Centros Universitários da região Sul tiveram dois cursos participantes que receberam os conceitos 3 e 4.
Foram sete cursos vinculados a Faculdades na região Sul, e nenhum destes ficou sem conceito.
Os demais receberam conceito 2 ( um curso ), conceito 3 ( cinco cursos, o conceito modal ) e conceito 5 ( também um curso ).
Na região Centro-Oeste, a metade dos seis cursos eram de Universidades.
Nesse tipo de organização, nenhum curso ficou sem conceito.
Os cursos receberam os conceitos 3 ( dois cursos ) e o conceito 4 ( um curso ).
Os Centros Universitários da região Centro-Oeste contaram com um curso, que recebeu o conceito 4.
Dos dois cursos em Faculdades na região Centro-Oeste, nenhum ficou sem conceito, um recebeu conceito 2, e o outro conceito 3.

% matched lemmas: questionário
\end{textsample}
